\chapter{49}\label{section-49}

\section{Auf der Alp, Bettmeralp/VS (CH), Mittwoch, 29.
August}\label{auf-der-alp-bettmeralpvs-ch-mittwoch-29.-august}

«Auf geht`s, ihr Lieben,» hörte Louis Joh in den weit offenen Raum
rufen. Sie hatten ihr sperriges Gepäck wieder an ihren Körpern
festgemacht und Steigeisen an den Bergschuhen befestigt. Die ersten
Meter gingen sie am Rande des Gletschers.

Nach einigen Schritten der Zunge entlang, würden sie das Eis überqueren.

Louis war aufgeregt wie lange nicht. Er blickte erwartungsvoll in die
Ferne auf die gegenüberliegende Seite der Gletscherzunge. Die steinerne
Seitenwand reichte hoch hinauf. Sie liess an mehreren Orten kleine
Plateaus erahnen.

Louis atmete tief ein. Die Luft erschien ihm schon fast unverschämt
rein. Dann blieb Joh stehen. Er zeigte mit dem Finger in Richtung
Felswand.

«Schau, Ciro,» er war ausser Atem, «dort oben sind sie. Siehst du die
winzig kleinen schwarz-weissen Punkte?»

Joh drehte sich zu Louis um. Der hob seine Sonnenbrille und suchte die
Wand mit den Augen ab.

«Nein, schau, mehr links und etwas höher...»

«Jetzt, ja,» Louis kniff die Augen zusammen, «und noch was Dunkles,
oder?»

«Genau, das sind die Schafe und dort geht gerade Mazi.»

Manuela war inzwischen an ihnen vorbei marschiert. Sie war schon am
Rande des Gletschers angekommen. Ihre Steigeisen hielt sie in den
Händen. Wieder stand sie mit weit ausgebreiteten Armen da. In ihrem
Gesicht konnte Louis ein feines Leuchten ausmachen. Wie oft musste sie
über die Jahre schon hier gestanden haben. Erstaunlich, wie sie sich
offenbar noch immer über diese Stimmung freuen konnte.

Bevor sie die letzten Meter in die Höhe in Angriff nahmen, blieben sie
stehen. Joh richtete mit ein paar Handgriffen Louis` Rucksack. Er
überprüfte die Stabilität des darüber geschnallten Gepäcks. Diesmal
kletterte Manuela dicht hinter Joh hinauf und Louis den beiden
hintennach.

«Warte, Joh,» rief Manuela eindringlich, «langsamer.»

Sie wirkte wacklig und stützte sich, zur Seite geneigt, mit einer Hand
an der Wand ab.

Louis starrte Manuela an. Erschrocken stellte er sich vor, wie sie
gleich die Balance verlieren und ihm entgegenschiessen würde. Er schloss
die Augen. Krallte sich an einem Felsvorsprung fest und dann stürzte das
Bild sich auf ihn. Es war, als sei die Szene direkt vor Louis vom Himmel
gefallen.

Dass Joh Manuela blitzschnell an der Hand hielt und die beiden
erleichtert lachten, nahm er nicht mehr wahr.

Er sah Giorgia vor sich. Er sah, wie sie schwankte. Sah, wie sie mit
diesem fremden Blick nach hinten verschwand.

Das Gewicht an seinem Körper schien ihn rücklings in den Abgrund zu
ziehen. Seine Hände zitterten. Schweiss tropfte ihm von der Stirn. Er
riss die Augen auf. Und schluckte.

«Ciro, alles ok?»

Johs Frage hallte von oben zu ihm hinunter. Louis schaute hoch. Er
staunte, wie weit er bereits zurücklag.

Sein Puls raste. Er konzentrierte sich auf seinen Atem. Ein und aus. Ein
und aus.

«Ja, ja, alles gut,» rief er zurück und war beruhigt, dass sie seinen
Schrecken nicht mitbekommen hatten.

«Nimm dir Zeit, keine Eile.»

Louis sah, wie Joh seinen Kopf zurückzog. Die beiden schienen oben
angekommen.

Er blieb starr stehen. Das explosive Gemisch aus Wut und Angst war
zurück. Noch vor wenigen Minuten hatte er sich gut gefühlt. Die Bilder
kamen hinterlistig und in den unmöglichsten Situationen. Er musste einen
Weg finden, sie im Zaun zu halten. Sie zu kontrollieren. Gerade hier,
umgeben von Felsen und steilen Abhängen hatten seine Erinnerungen
leichtes Spiel. Sie mischten sich gewaltsam ein, erschreckten und
lähmten ihn. Er musste sich eine Strategie ausdenken. Bevor Schlimmeres
geschah.
